\section{An explicit order bound in the qualified transitive-core gate}
\label{sec:explicit-inducing-order}

\paragraph{Scope and status.}
HD-FQ01--03 is a new explicit quantitative consequence of the accepted
family reductions, HD-TW03 and HD-TC01--03.  The complete author and root whole-proof readings pass. Root also read
the entire author checker, its first complete output, and the entire
source-and-status record; no fresh root execution or fresh full
reading of all previously accepted source interfaces is claimed.
This explicit quantitative synthesis is recorded as derived work
for the formal movement rule.
Fix one of the eight non-Eisenstein class-number-one imaginary quadratic
fields
\[
 K\in\{\mathbf Q(i),\mathbf Q(\sqrt{-2}),\mathbf Q(\sqrt{-7}),
 \mathbf Q(\sqrt{-11}),\mathbf Q(\sqrt{-19}),
 \mathbf Q(\sqrt{-43}),\mathbf Q(\sqrt{-67}),
 \mathbf Q(\sqrt{-163})\}.
 \tag{FQ.1}
\]
Consider an actual finite-cover Weil sixfold source of exact original
character field $K$, period rank nine, and nonsplit actual polarization.
Let $\tau$ be an absolutely irreducible inducing block of degree $e\ge10$,
and let $J$ be its finite faithful matrix image.  Assume $J$ is primitive
and tensor indecomposable in the \emph{factorwise} sense of D-BI-19.
This convention permits permutations of tensor factors.  It does not
permit a nontrivial tensor decomposition preserved factorwise by $J$.
No equality of the original and inducing character fields is assumed.

These are the precise hypotheses throughout this section.  In particular,
no result here bounds inducing blocks of degree below ten, covers the
remaining factorwise tensor-decomposable case, or bounds an arbitrarily
larger original group from which $\tau$ is induced.

\subsection*{HD-FQ01. A common simple-quotient order ceiling}

Define the following positive integers, without expanding their decimal
representations:
\[
 \begin{split}
 M_*&=1+241\cdot240!\cdot2^{239},\qquad F_*=240M_*+2,\\
 \mathcal N&=[100(1+M_*)F_*^{240}]^2,\qquad D_0=10^{62},\\
 S_0&=\max\{\mathcal N!,\ 27^{3D_0^2},\ 10^{60}\}.
 \end{split}                                                   \tag{FQ.2}
\]

\begin{lemma}[HD-FQ01]
Every simple quotient of a quasisimple factor of the D-BI-19 normal core
has order at most $S_0$.
\end{lemma}

\begin{proof}
D-BI-19 supplies a normal irreducibly acting central product
$E_1\cdots E_t$, with the quasisimple factors conjugate under $J$.
Write $S=E_i/Z(E_i)$ for their common nonabelian simple quotient.
We use the finite simple group classification in exactly its usual
alternating, Lie-type and sporadic form; a complete statement with the
classical and exceptional families is given in
\cite[\S2.2.1, author p.~13]{HES:LiXia2022}.
This proof does not reconstruct that classification.

If $S$ is alternating, the ordinary representations with $n\ge\mathcal N$
are excluded by HD-TW03 in the present non-Eisenstein nonsplit scope.
The genuine spin alternatives for $n\ge8$ are excluded by HD-SPIN-02;
the exceptional covers in smaller alternating degrees do not affect
the inequality $n<\mathcal N$.  Thus
\[
                         |S|=n!/2\le\mathcal N!.
                                                               \tag{FQ.3}
\]
The spin and ordinary alternatives here concern the actual represented
central quotient; no ordinary-character assumption is imposed on $E_i$.

If $S$ is not alternating and is of Lie type over $\mathbf F_q$,
D-BI-20 applies under the present hypotheses and gives $q\le27$.
The accepted large-rank bounds then give the following deliberately
coarse ceilings for its classical natural dimension $d$:
\[
\begin{array}{c|c|l}
\text{family}&\text{ceiling sufficient here}&\text{accepted interface}\\ \hline
\operatorname{PSL}_{d}(q)&d<10^{61}&\text{D-LNW-01 and D-BI-23}\\
\operatorname{PSU}_{d}(q)&d<10^{61}&\text{D-UNW-03 and D-BI-24}\\
\operatorname{PSp}_{d}(q)&d\le6&\text{D-HES-01--06 and D-HOS-03--22}\\
\text{orthogonal, }q\text{ even}&d\le1134&\text{D-OL-11}\\
\text{orthogonal, }q\text{ odd}&d\le1134&\text{D-OL-19--24}.
\end{array}                                                    \tag{FQ.4}
\]
The linear and unitary statements include their level-one alternatives,
as specified in the cited interfaces.  The symplectic statements cover
all represented central kernels in the large ranks; the exceptional
small-rank multipliers are left inside the displayed finite range.
In even characteristic the large-rank orthogonal simply connected cover
has trivial center and no exceptional multiplier, as in the multiplier
scope already used in D-OL-11.  In odd characteristic D-OL-19--24 include
both ordinary and genuinely spin representations and arbitrary
represented central kernels.  In particular the withdrawn stronger
threshold based on an incorrectly read second square root in LST is
\emph{not} an input.  Every entry in (FQ.4) is less than $D_0$.

For these classical groups their defining special linear, unitary,
symplectic or orthogonal matrix group surjects onto the simple group.
Its natural degree is $d$, and its matrix coefficients lie in
$\mathbf F_q$ or $\mathbf F_{q^2}$.  Counting all matrices gives
$|S|<q^{2d^2}$, which is less than $27^{3D_0^2}$.
For clarity, this estimate does not require the exact classical order
formula or a faithful linear model for the simple quotient itself.

For exceptional groups use the published matrix constructions in
\cite[\S3, author pp.~4--12; \S4, pp.~12--17]{FQ:HRT2001}.
The untwisted degrees are $7,26,27,56,248$ for
$G_2,F_4,E_6,E_7,E_8$.  The twisted constructions have degrees
$4,8,27,26,7$ for
${}^2B_2,{}^3D_4,{}^2E_6,{}^2F_4,{}^2G_2$, over extensions of
$\mathbf F_q$ of degree at most three.  Their matrix groups map
onto the simple group, or contain its Tits derived subgroup in the
stated ${}^2F_4(2)$ case.  Hence every exceptional simple group here has
\[
                        |S|<q^{3\cdot248^2}
                                 \le27^{3D_0^2}.
                                                               \tag{FQ.5}
\]
The quoted constructions include the bad-characteristic parameters;
irreducibility of the displayed matrix module is unnecessary for
this order estimate.  The defining-field convention is $q$, not the
larger matrix coefficient field.

Finally all twenty-six sporadic orders are below $10^{60}$.
The classification list and the range up to the Monster, whose order
is less than $10^{54}$, are explicitly recorded in
\cite[Lecture~1, p.~1]{FQ:Wilson2008}; the exact Monster order is also
listed in the primary ATLAS data \cite{FQ:AtlasMonster}.
This harmless envelope does not require reopening the earlier stronger
sporadic exclusion.  Combining these cases proves the lemma.
\end{proof}

\subsection*{HD-FQ02. Bounding every quasisimple central cover}

\begin{lemma}[HD-FQ02]
For a finite group $S$ of order $s$, its integral Schur multiplier satisfies
\[
                  |H_2(S,\mathbf Z)|\le4^{(s-1)^2}.
                                                               \tag{FQ.6}
\]
Consequently, if $S$ is nonabelian simple and $E\to S$ is any
quasisimple central extension, then
\[
                         |E|\le s\,4^{(s-1)^2}.
                                                               \tag{FQ.7}
\]
\end{lemma}

\begin{proof}
Use the normalized integral bar resolution, followed by trivial
coinvariants.  Its free abelian groups in degrees two and three have
bases $[g|h]$ and $[g|h|k]$ with each entry nonidentity; any symbol with
an identity entry is interpreted as zero.  The differential is
\[
 d_3[g|h|k]=[h|k]-[gh|k]+[g|hk]-[g|h].
                                                               \tag{FQ.8}
\]
The complete resolution and its use to compute integral group homology
are given in \cite[Definition~6.5.1, Theorem~6.5.3 and
Application~6.5.4, pp.~177--179]{FQ:Weibel1994}.
Each column of this matrix is a sum of four signed coordinate vectors,
so its Euclidean norm is at most four even when some terms coincide.
Its row count is $(s-1)^2$.

The chain-complex identity gives the exact sequence
\[
 0\longrightarrow H_2(S,\mathbf Z)
  \longrightarrow\operatorname{coker}(d_3)
  \longrightarrow\operatorname{im}(d_2)\longrightarrow0.
                                                               \tag{FQ.9}
\]
The last group is free abelian, being a subgroup of a free abelian
group.  The first is finite: integral positive-degree homology of a
finite group is finitely generated and killed by $s$, by
\cite[Theorem~6.5.8 and Corollaries~6.5.9--6.5.10, p.~180]{FQ:Weibel1994}.
It follows that $H_2(S,\mathbf Z)$ is exactly the torsion subgroup of
$\operatorname{coker}(d_3)$.

Let $r$ be the rational rank of $d_3$.  Smith normal form says that
the order of that torsion subgroup is the greatest common divisor
of the nonzero $r$-rowed minors.  Choose any one such minor.
Deleting rows cannot increase column norm, so Hadamard's determinant
inequality bounds its absolute determinant by $4^r$.
Since $r\le(s-1)^2$, this proves (FQ.6).
For $r=0$ the torsion group is trivial and the same conclusion holds.

A perfect group has a universal central extension with kernel
$H_2(S,\mathbf Z)$; see the complete construction and proof in
\cite[\S6.9, Theorem~6.9.5, pp.~198--201]{FQ:Weibel1994}.
Its map to any central extension $E\to S$ is surjective when $E$
is perfect: its image $U$ maps onto $S$, so $E=U Z$ for the
central kernel $Z$, and then $E=[E,E]=[U,U]\subset U$.
Every quasisimple $E$ is therefore a quotient of that universal
extension.  Equation (FQ.7) follows from (FQ.6).
No classification of Schur multipliers is used in this new estimate.
\end{proof}

\subsection*{HD-FQ03. Explicit inducing order and its arithmetic consequence}

Set
\[
 \begin{split}
 A_0&=S_0\,4^{S_0^2},\qquad T_0=2A_0,\qquad E_0=A_0^{T_0},\\
 B_0&=S_0!,\qquad
 Q_0=E_0\bigl(B_0^{T_0}T_0!\bigr)^2.
 \end{split}                                                   \tag{FQ.10}
\]
All these are explicitly specified integers; no computational
enumeration up to them is asserted.

\begin{theorem}[HD-FQ03]
Every source under the initial hypotheses has
\[
                 t\le T_0,\qquad e\le E_0,\qquad |J|\le Q_0.
                                                               \tag{FQ.11}
\]
Its connected inducing $J$-cover $Y\to X=\mathbf P^1$ has
\[
                  r\le2E_0+6,\qquad
                  g(Y)\le1+Q_0(E_0+2).
                                                               \tag{FQ.12}
\]
For each fixed $K$ in (FQ.1), set
\[
                M_0=\operatorname{lcm}(8,|\operatorname{disc}K|,
                                      1,2,\ldots,Q_0).
                                                               \tag{FQ.13}
\]
Every prime $p\equiv-1\pmod{M_0}$ is excluded as the signed
determinant class of a naturally polarized source in this qualified
class.  The same exclusion holds for a target with
$\operatorname{End}^0=K$ which is $K$-isogenous to such a source.
There are infinitely many pairwise distinct excluded norm classes.
\end{theorem}

\begin{proof}
By HD-FQ01--02 each actual quasisimple factor has order at most $A_0$.
An irreducible complex representation of a finite group has degree
at most its group order, so its factor degree satisfies $2\le a\le A_0$.
D-BI-19 gives equal factor degrees and $e=a^t$.
HD-TC02 gives
\[
 t\le\max\{4,\lceil5a/4\rceil-1\}\le2A_0=T_0,
                 \qquad e\le A_0^{T_0}=E_0.
                                                               \tag{FQ.14}
\]
An automorphism of a finite group of order $s$ permutes its underlying
set faithfully, so $|\operatorname{Aut}(S)|\le s!\le B_0$.
For $P=S^t$, the simple direct factors are the minimal normal subgroups;
therefore every automorphism permutes them and restricts to an
automorphism on each.  Thus
\[
 |\operatorname{Aut}(P)|=|\operatorname{Aut}(S)|^t t!
                                  \le B_0^{T_0}T_0!.
                                                               \tag{FQ.15}
\]
HD-TC03, which retains the entire scalar kernel, now yields
$|J|\le e|\operatorname{Aut}(P)|^2\le Q_0$.
That theorem uses the actual generating inertia tuple: all but possibly
one inertia have a fixed vector, their projective powers control their
determinants, and product one controls the last determinant as well.
Consequently the estimate does not assume a bounded scalar group
before using the source geometry.
The same source Euler identity and Galois Riemann--Hurwitz argument
in HD-TC03 give (FQ.12).

Apply HD-PD02 with its bounded inducing order $Q=Q_0$.
For every prime in (FQ.13), that theorem proves that the natural
rank-six signed Hermitian determinant has even $p$-valuation.
The local proof allows $p$ to divide the order of the original group
or of the discarded kernel; only $p\nmid|J|$ is needed, and the
congruence implies it.  Rank-six polarization rescaling contributes
a sixth power, which is a norm, so the target comparison under
$\operatorname{End}^0=K$ is exactly the one stated in HD-PD02.
The signed class $p$ has odd valuation and is therefore impossible.
Dirichlet's theorem gives infinitely many such primes.  They are inert
in $K$, and distinct such primes give different classes in
$\mathbf Q^\times/\operatorname{Nm}_{K/\mathbf Q}(K^\times)$,
since norms have even valuation at inert primes.
These are actual signature-$(3,3)$ target types, represented by
$\operatorname{diag}(1,1,1,-1,-1,-p)$.
\end{proof}

\paragraph{What the explicit bound decides.}
The multiplication-table, character, tuple and complex period-rank
enumeration in the refined HD-TC03 is now finite for the qualified
class above, with the displayed uniform integer ceilings.
Its original-$K$ descent condition remains separate, and no such
enormous enumeration has been performed.
D-OG01--02 exhibit why this cannot be promoted to a bound on the
original group or original induction index: arbitrary wreath padding
can preserve an actual primitive sixfold factor and its determinant.
The arithmetic conclusion excludes this specified source route for
infinitely many target determinant classes; it is not a claim that
the associated Hodge classes are nonalgebraic.
The remaining factorwise tensor-decomposable case and degrees below
ten are outside this theorem, so it supplies no unconditional
all-supplier obstruction.
